% Problem formulation (draft, Week 11; written 4 Oct 2026).
% Notation follows docs/v2/SIM_SPEC.md and PROTOCOL.md; numbers will come from numbers.tex.

\section{Problem formulation}
\label{sec:problem}

A cluster of hosts runs virtual machines (VMs). A host fails when any of its components fails
(drives, or other parts), and every VM on a host that fails is interrupted. Once a day an
operator may live-migrate some VMs to other hosts~\cite{clark2005live}; a migration consumes
network and host resources and causes a short blackout, so migrations are limited by a daily
budget. The question is which VMs to move.

\begin{table}[t]
  \centering
  \small
  \begin{tabular}{@{}ll@{}}
    \toprule
    Symbol & Meaning \\
    \midrule
    $j$, $i$ & host; VM (on host $j(i)$) \\
    $H$ & risk horizon in days (30 unless stated) \\
    $p_j$ & probability that host $j$ fails within $H$ days, given it is up now \\
    $c_i$ & cost if VM $i$ is interrupted: tier weight $\times$ recovery downtime, weighted by the
            probability that $i$ is still running at the failure \\
    $e_i$ & probability that $i$'s migration completes before its host fails (depends on $i$'s
            position in the host's migration queue) \\
    $m_i$ & cost of migrating $i$ (blackout and resource time, in the same units as $c_i$) \\
    $w_i$ & resource $i$'s migration consumes from the budget (memory, GB) \\
    $B$ & daily budget, in GB migrated per cluster \\
    $x_i$ & decision: migrate $i$ today ($x_i \in \{0,1\}$) \\
    \bottomrule
  \end{tabular}
  \caption{Notation.}
  \label{tab:notation}
\end{table}

\paragraph{Decision problem.}
Migrating VM $i$ changes the expected cost by the benefit
\begin{equation}
  b_i = p_{j(i)}\, c_i\, e_i - m_i ,
  \label{eq:benefit}
\end{equation}
the expected interruption cost avoided minus the cost of moving. Each day the operator solves
\begin{equation}
  \max_{x}\ \sum_i b_i x_i
  \quad \text{subject to} \quad \sum_i w_i x_i \le B,\quad x_i \in \{0,1\}.
  \label{eq:knapsack}
\end{equation}
This is a 0--1 knapsack. Its fractional relaxation is solved exactly by ranking VMs by the
density $b_i / w_i$ and filling the budget while $b_i > 0$; the integral problem differs from it
by at most one item. Because $e_i$ depends on the position of $i$ in its host's queue, the plan
is built greedily with $e$ recomputed after each pick. We call this rule \emph{budget-aware
migration} (BAM). Cost-sensitive decision rules of this form are classical~\cite{elkan2001foundations};
the questions this paper asks are when ranking by density beats the operational default of
ranking by $p_j$ alone, and by how much.

\paragraph{Baseline.}
The threshold policy migrates VMs on hosts with $p_j \ge \tau$. At an equal budget it is the
special case of~\eqref{eq:knapsack} that ranks by $p_j$ and ignores $c$, $m$, $e$ and $w$. We
tune $\tau$ and the order in which a host's VMs are moved on held-out seeds, and give it exactly
the same budget in GB as every other policy.

\paragraph{Objective and accounting.}
Policies are compared by the cost that actually occurs in simulation: the interruption cost of
every VM that is on a host when it fails, plus the cost of every migration. No migration is
credited with an averted failure; an interruption that does not happen simply does not appear.
The horizon $H$ of the probability a policy uses is the horizon over which its consequences are
observed. Our earlier course prototype credited migrations with failures months after a 30-day
decision; that accounting, not the decision rule, explained most of its apparent loss to the
threshold at zero correlation (Section~\ref{sec:v1}).

\paragraph{When can density ranking help?}
If $c_i e_i / w_i$ were the same for every VM, both rules would pick the same VMs. Density
ranking can only gain when this ratio varies across VMs, and gains most when the VMs with a
high ratio sit on the riskiest hosts. We therefore describe a cluster by the correlation between
VM criticality and host risk, $\rho$, and by the spread of $\log(c e / w)$, and measure where real
placement lies on that axis (Section~\ref{sec:rho}).
